% Figura: mapa del capítulo sobre el pipeline.
% Se usa en la apertura del Capítulo 2. Cada componente lleva la sección que lo explica. Es la versión
% simplificada del diagrama del pipeline de sec:met_pipeline: si uno cambia, revisar el otro.

\begin{figure}[htbp]
    \centering
    \begin{tikzpicture}[x=1mm, y=1mm]

        \node[figentrada, text width=26mm] (img) at (-4, 0) {Imagen\\dermatoscópica};

        \node[figtarea, text width=30mm] (clf) at (40, 13)
            {Clasificador\\{\footnotesize visión, §\ref{sec:marco_vision}}};
        \node[figtarea, text width=30mm] (knn) at (40, -13)
            {Codificador visual y kNN\\{\footnotesize §\ref{sec:marco_vlm}}};

        \node[figtarea, text width=36mm] (gen) at (86, 0)
            {Generador\\{\footnotesize lenguaje, §\ref{sec:marco_llm};\\visión--lenguaje, §\ref{sec:marco_vlm}}};

        \node[figsalida, text width=18mm] (rep) at (126, 0) {Reporte};

        \node[fignodo, text width=30mm, draw=figpurpura!70!black, fill=figpurpura!12] (eval)
            at (126, -24) {Evaluación, §\ref{sec:marco_evaluacion}\\{\footnotesize generalización,
            §\ref{sec:marco_generalizacion}}};

        \draw[figflecha] (img.east) -- ++(5, 0) |- (clf.west);
        \draw[figflecha] (img.east) -- ++(5, 0) |- (knn.west);
        \draw[figflecha] (clf.east) -- node[figetiqueta, above] {categoría} ++(12, 0) |- ([yshift=3mm]gen.west);
        \draw[figflecha] (knn.east) -- node[figetiqueta, below, align=center] {casos\\parecidos} ++(12, 0) |- ([yshift=-3mm]gen.west);
        \draw[figflecha] (gen.east) -- (rep.west);
        \draw[figflecha] (rep.south) -- (eval.north);

        \begin{scope}[on background layer]
            \node[figgrupo={todos los componentes: fundamentos, §\ref{sec:marco_fundamentos};
                  atención, §\ref{sec:marco_atencion}}, fit=(clf) (knn) (gen)] {};
        \end{scope}

    \end{tikzpicture}
    \caption[El capítulo recorrido sobre el \emph{pipeline}]{El capítulo recorrido sobre el
    \emph{pipeline}. Cada componente lleva la sección que presenta los conceptos que usa. El clasificador
    asigna una categoría a la imagen, un codificador visual recupera casos parecidos y el generador redacta
    el reporte con esa información; después la salida se evalúa y se examina si lo medido se sostiene
    fuera de los datos de entrenamiento. Los fundamentos del entrenamiento y la operación de atención son
    comunes a los tres componentes.}
    \label{fig:marco_mapa}
\end{figure}
